%~~~~~~~~~~~~~~~~~~~~~~~~~~~~~~~~~~~~~~~~~~~~~~~~~~~~~~~~~~~~~~~~~~~~~~~~~~~~~~~~~~~~~~~~~%
%	                                     ABSTRACT   	                                  %
%~~~~~~~~~~~~~~~~~~~~~~~~~~~~~~~~~~~~~~~~~~~~~~~~~~~~~~~~~~~~~~~~~~~~~~~~~~~~~~~~~~~~~~~~~%
\begin{flushleft} % This section is not essential for the abstract
    \setlength{\parskip}{0pt}
        %\bigskip
    {\centering{{\Large{\bf{ABSTRACT}}}} \par}
    \medskip
    \vspace{6pt}
    \hrule \vspace{0.4cm}
		  Student Name:	 \textbf{\ThesisAuthor}	\hfill Reg. No.: \textbf{\RegNumber} \\
		Degree:  \textbf{\DegreeDisplay} \hfill Dept.:  \textbf{School of \Department} \\
		Thesis Title: \textbf{\ThesisTitle}\\
		Thesis Supervisor:  \textbf{\ProjectSupervisor}\\
		Date of thesis submission: \textbf{{ \today}\\[0.4cm] }
	\hrule \vspace{1.5cm} 
\end{flushleft}
\noindent
Numerical Relativity can be defined as the study of understanding GR numerically in such a way that the both of their solutions agree with each other at the limit of infinite resolution. This has open doors to understanding relativistic systems which were hard or impossible to attempt to using traditional theoretical tools. In this regard, one of the most important insights regarding GR that came about purely due to NR is the discovery of Critical Phenomenon in gravitational collapse.

In this 
\vspace{4cm}

{\large \textbf{Keywords: }}\par{\Large \ThesisKeywords}

\cleardoublepage